\SectionStart
\section{问题二：模型建立与求解}
\subsection{目标与上下游接口}
\WritingContract{明确本问要回答的量化问题及其与前后问的连接。}{写明本问输入、前问传入量、输出、决策或状态变量、主要约束及下游接口。}{与题面、question manifest 和程序接口一致，传递量注明单位和粒度。}{不得用“基于问题一”替代具体接口，不得提前给结果。}
\subsection{数据特征或机理依据}
\WritingContract{说明模型结构由哪些数据现象、前问结果或领域关系推出。}{给出影响建模的分布、趋势、相关结构、守恒关系或机理，并解释前问输出的作用。}{只使用可定位的数据剖析、冻结上游结果、题面关系或已核验来源。}{不得堆放无关统计，不得把相关性直接解释为因果。}
\subsection{模型选择与备选方案比较}
\WritingContract{论证当前模型如何处理本问的关键冲突。}{比较主模型、同输出基线及至少一个备选思路，说明选择标准、适用条件和不优先原因。}{依据数据结构、约束性质、可解释性、计算规模或真实文献。}{不得固定套用某类算法，不得用未经实验的性能判断支持选择。}
\subsection{模型建立}
\WritingContract{把本问要求及上游接口转化为可计算模型。}{定义新增变量和参数，给出目标、约束、边界、单位、跨问连接及必要的线性化或近似。}{公式与题目变量、前问输出和代码实现一致，参数有来源。}{不得照抄通式，不得忽略接口单位、变量域或边界条件。}
\subsection{求解算法}
\WritingContract{说明模型如何被稳定、可复现地求解。}{写明输入、初始化、关键步骤、求解器或算法参数、随机种子、停止条件、运行规模和硬约束复核。}{步骤对应正式 runner 和运行清单，单问与全流程调用结果一致。}{不得粘贴大段工程代码，不得声称未实现的算法、复杂度或最优性。}
\subsection{核心结果与解释}
\WritingContract{通过统一输出比较直接回答本问。}{报告冻结核心数值、单位和基线，展示必要的结构、分布或权衡，并按“结果—原因—意义”解释。}{主模型与基线输出类别一致，表图和文字引用同一冻结证据。}{不得手填数字，不得只描述图形外观，不得隐去不利但重要的权衡。}
\subsection{模型检验}
\WritingContract{验证结果不是由偶然参数、不可行解或单一场景造成。}{依据题型给出误差、可行性、守恒、灵敏度、压力场景、多种子或替代方法比较。}{报告检验范围、样本或场景、指标和是否改变结论。}{不得以单次运行、随机结果或收敛图替代基线与稳健性证据。}
\subsection{本问结论与适用边界}
\WritingContract{收束答案并说明对后续问题的可用输出。}{概括最终方案或判断、相对基线的结论、验证结果、成立条件、失败边界和下游接口。}{只引用本问及冻结上游证据中的正式主张。}{不得新增结果，不得把局部或近似求解宣称为全局普适最优。}
